\begin{tikzpicture}
    % Un componente riusabile della figura: pics/<nome>/.style n args={3}.
    % Argomenti: #1 = nome (serve a dare nomi unici ai nodi), #2 = posizione, #3 = intervalli di overlay in cui
    % il gruppo è evidenziato (anche non contigui). Si chiama con \pic{box group={nome}{(x,y)}{intervalli}};
    \tikzset{
        pics/box group/.style n args={3}{
            code={
                \begin{scope}[shift={#2}]
                    \foreach \i in {1,2,3}{
                        \node[rboxes, minimum width=4em,
                            on slide={<#3> highlight light}] (box-#1-\i) at (0,{-0.9*\i}) {#1.\i};
                    }
                \end{scope}
            }
        }
    }

    \pic{box group={a}{(0,0)}{2-2, 4-4}};
    \pic{box group={b}{(5,0)}{3-4}};

    % I nodi creati dal componente si usano come qualunque altro nodo: box-<nome>-<numero>.
    \draw[->, link, on slide={<4-4> highlight link}] (box-a-3) -- node[on link] {relation} (box-b-3);
\end{tikzpicture}
